% Part: computability
% Chapter: machines-computations
% Section: halting-problem

\documentclass[../../../include/open-logic-section]{subfiles}

\begin{document}

\olfileid{tur}{und}{hal}
\olsection{থামার সমস্যা}

\begin{explain}
ধরা যাক, টুরিং যন্ত্রের বর্ণনাগুলির কোনো একটি তালিকায়ন আমরা স্থির করেছি।
তাহলে প্রতিটি টুরিং যন্ত্র একটি \emph{সূচক} পায়: টুরিং-যন্ত্রের বর্ণনার
তালিকা $M_1$, $M_2$, $M_3$, \dots{}-তে তার স্থান।

আমরা জানি, টুরিং-অগণনসাধ্য অপেক্ষক অবশ্যই আছে: টুরিং যন্ত্রের
বর্ণনাগুলির সেট—এবং তাই টুরিং যন্ত্রের সেট—!!{enumerable}, কিন্তু
$\Nat$ থেকে $\Nat$-এ সকল অপেক্ষকের সেট তা নয়। তবে অগণনসাধ্য অপেক্ষকের
নির্দিষ্ট উদাহরণও পাওয়া যায়। তেমন একটি অপেক্ষক হলো থামা-অপেক্ষক।
\end{explain}

\begin{defn}[থামা-অপেক্ষক]
 \emph{থামা-অপেক্ষক}~$h$ নিচের মতো সংজ্ঞায়িত:
\[
h(e,n) =
\begin{cases}
  \text{0} & \text{যদি যন্ত্র~$M_e$ নিবেশ $n$-এ না থামে} \\
  \text{1} & \text{যদি যন্ত্র~$M_e$ নিবেশ $n$-এ থামে}
\end{cases}
\]
\end{defn}

\begin{defn}[থামার সমস্যা]
\emph{থামার সমস্যা} হলো, যেকোনো $e$, $n$-এর জন্য টুরিং যন্ত্র~$M_e$
~$n$টি দাগের নিবেশে থামে কি না তা নির্ধারণ করার সমস্যা।
\end{defn}

\begin{explain}
একটি সম্পর্কিত অপেক্ষক~$s$ টুরিং-গণনসাধ্য নয় দেখিয়ে আমরা দেখাব
যে~$h$ টুরিং-গণনসাধ্য নয়। প্রমাণটি এই দুই তথ্যের উপর নির্ভর করে: টুরিং
যন্ত্রে যা কিছু গণনা করা যায়, শৃঙ্খলাবদ্ধ টুরিং যন্ত্রেও তা গণনা করা যায়
(\olref[mac][dis]{sec}); এবং দুটি টুরিং যন্ত্রকে জুড়ে একটিমাত্র যন্ত্র
তৈরি করা যায় (\olref[mac][cmb]{sec})।
\end{explain}

\begin{defn} অপেক্ষক~$s$ নিচের মতো সংজ্ঞায়িত:
\[
s(e) =
\begin{cases}
  \text{0} & \text{যদি যন্ত্র~$M_e$ নিবেশ $e$-তে না থামে} \\
  \text{1} & \text{যদি যন্ত্র~$M_e$ নিবেশ $e$-তে থামে}
\end{cases}
\]
\end{defn}

\begin{lem}
অপেক্ষক~$s$ টুরিং-গণনসাধ্য নয়।
\end{lem}

\begin{proof}
স্ববিরোধ পাওয়ার উদ্দেশ্যে ধরি, অপেক্ষক~$s$ টুরিং-গণনসাধ্য। তাহলে~$s$
গণনাকারী একটি টুরিং যন্ত্র~$S$ থাকত। সাধারণত্ব ক্ষুণ্ণ না করে ধরে নিতে
পারি,~$S$ থামলে প্রথম ঘরটি পড়তে পড়তেই থামে (অর্থাৎ সেটি শৃঙ্খলাবদ্ধ)।
এই যন্ত্রটিকে অন্য একটি যন্ত্র~$J$-এর সঙ্গে “জোড়া” যায়। যন্ত্রটি নিবেশ~$0$-এ
শুরু করলে (অর্থাৎ আরম্ভিক দশায় ফিতার শেষ-চিহ্নের ডান দিকের ঘরটি পড়ার
সময়~$\TMblank$ পেলে) থামে; অন্যথায় ডান দিকে চলে যায় এবং কখনও থামে না।
$S \frown J$, অর্থাৎ~$S$-কে~$J$-এর সঙ্গে জুড়ে তৈরি যন্ত্রটি একটি টুরিং
যন্ত্র; তাই কোনো~$e$-এর জন্য সেটি~$M_e$ (অর্থাৎ তালিকায় কোথাও সেটি আছে)।
~$M_e$-কে $e$টি~$\TMstroke$-এর নিবেশে শুরু করো। দুটি সম্ভাবনা আছে:
হয়~$M_e$ থামে, নয়তো থামে না।
\begin{enumerate}
\item ধরা যাক, $e$টি~$\TMstroke$-এর নিবেশে~$M_e$ থামে। তাহলে $s(e)
  = 1$। তাই~$S$ নিবেশ~$e$-তে শুরু করলে ফিতায় নির্গম হিসেবে একটিমাত্র
  ~$\TMstroke$ রেখে থামে। এরপর~$J$ ফিতায় একটি~$\TMstroke$ নিয়ে শুরু করে।
  সেক্ষেত্রে~$J$ থামে না। কিন্তু~$M_e$ হলো যন্ত্র $S \frown J$; তাই
  ~$S$-এর পরে~$J$ যা করত, সেটিরও ঠিক তাই করার কথা (অর্থাৎ এই ক্ষেত্রে
  ডান দিকে চলে গিয়ে কখনও না-থামা)। অতএব $e$টি~$\TMstroke$-এর নিবেশে
  ~ $M_e$ থামতে পারে না।

\item এবার ধরা যাক, $e$টি~$\TMstroke$-এর নিবেশে~$M_e$ থামে না। তাহলে
  $s(e) = 0$, এবং~$S$ নিবেশ~$e$-তে শুরু করে ফাঁকা ফিতা নিয়ে থামে।~$J$
  ফাঁকা ফিতায় শুরু করলে সঙ্গে সঙ্গে থামে। আবার~$M_e$ ঠিক~$S$-এর পরে~$J$
  যা করত তাই করে; ফলে $e$টি~$\TMstroke$-এর নিবেশে~$M_e$ অবশ্যই থামে।
\end{enumerate}
প্রতিটি ক্ষেত্রেই আমরা আমাদের অনুমানের সঙ্গে স্ববিরোধে পৌঁছাই। এতে প্রমাণ
হয়, এমন কোনো টুরিং যন্ত্র~$S$ থাকতে পারে না:~$s$ টুরিং-গণনসাধ্য নয়।
% BN-SRC-177 TARGET-CORRECTION-BEGIN
\par\smallskip\noindent\textnormal{\small\textbf{উৎস-সংশোধন (BN-SRC-177):} হিমায়িত ইংরেজি প্রমাণে যন্ত্র-সংযোজনকে দুবার $S \concat J$ লেখা হয়েছে, যদিও আগের সংজ্ঞায় $M \frown M'$ হলো যন্ত্র-সংযোজন আর $\concat$ প্রতীকক্রমের সংযুক্তকরণ। বাংলা পাঠে উভয় যন্ত্র-সংযোজনে $S \frown J$ পুনরুদ্ধার করা হয়েছে।} % BN-SRC-177 TARGET-CORRECTION-END
\end{proof}

\begin{thm}[থামার সমস্যার অসমাধানযোগ্যতা]
\ollabel{thm:halting-problem} থামার সমস্যা অসমাধানযোগ্য, অর্থাৎ অপেক্ষক
~$h$ টুরিং-গণনসাধ্য নয়।
\end{thm}

\begin{proof}
ধরা যাক, $h$ টুরিং-গণনসাধ্য—ধরা যাক, কোনো টুরিং যন্ত্র~$H$ সেটি গণনা
করে।~$H$ ব্যবহার করে~$s$ গণনাকারী একটি টুরিং যন্ত্র তৈরি করতে পারতাম:
প্রথমে নিবেশটির একটি অনুলিপি তৈরি করো (দুটি অনুলিপির মাঝে একটি~$\TMblank$
প্রতীক রেখে)। তারপর আবার শুরুতে ফিরে~$H$ চালাও। আগের কাজটি করার যন্ত্র যে
স্পষ্টতই তৈরি করা যায়, তা আমরা জানি (\cref{tur:mac:dis:prob:copier}); আর
$H$ থাকলে এমন অনুলিপিকারী যন্ত্রের সঙ্গে সেটিকে “জুড়ে” নতুন একটি যন্ত্র
পেতাম, যা~$M_e$ নিবেশ~$e$-তে থামে কি না নির্ধারণ করে, অর্থাৎ~$s$ গণনা
করে। কিন্তু আমরা ইতিমধ্যেই দেখিয়েছি, এমন কোনো যন্ত্র থাকতে পারে না। তাই
~$h$-ও টুরিং-গণনসাধ্য নয়।
\end{proof}

\begin{prob}
তিন-থামা (3-Halt) সমস্যা হলো একটি নির্ণয়-পদ্ধতি দেওয়ার সমস্যা, যা
নির্ধারণ করবে যেকোনো নির্বাচিত টুরিং যন্ত্র অন্যথায় ফাঁকা ফিতায় তিনটি
~$\TMstroke$-এর নিবেশে থামে কি না। প্রমাণ করো যে তিন-থামা সমস্যা
অসমাধানযোগ্য।
\end{prob}

\begin{prob}
দেখাও, যেসব টুরিং-যন্ত্র ও নিবেশ যুগল $M_e$ ও~$n$-এর জন্য $e \neq n$,
তাদের ক্ষেত্রে থামার সমস্যা সমাধানযোগ্য হলে $e = n$ হওয়ার ক্ষেত্রগুলিতেও
সেটি সমাধানযোগ্য।
\end{prob}

\begin{prob}
আমরা প্রমাণ করেছি, নিবেশটি যদি সংখ্যা~$e$ হয় এবং সব টুরিং যন্ত্রের একটি
তালিকায়নের মাধ্যমে সেটি টুরিং যন্ত্র~$M_e$-কে শনাক্ত করে, তবে থামার সমস্যা
অসমাধানযোগ্য। তার বদলে~\olref[tur][und][enu]{sec}-এর টুরিং যন্ত্রের
বর্ণনাগুলিকেই যদি সরাসরি নিবেশ হিসেবে অনুমোদন করি, তবে কী হয়? সূচকের
বদলে টুরিং যন্ত্রের বর্ণনা নিবেশ হিসেবে নেয় অথচ থামার সমস্যা নির্ণয়
করে—এমন কোনো টুরিং যন্ত্র কি থাকতে পারে? কেন পারে অথবা পারে না, ব্যাখ্যা
করো।
\end{prob}

\begin{prob} দেখাও যে \emph{আংশিক} অপেক্ষক~$s'$ নিচের মতো সংজ্ঞায়িত হলে
  \[
  s'(e) =
  \begin{cases}
    \text{1} & \text{যদি যন্ত্র~$M_e$ নিবেশ $e$-তে থামে}\\
    \text{অসংজ্ঞায়িত} & \text{যদি যন্ত্র~$M_e$ নিবেশ $e$-তে না থামে}
  \end{cases}
  \]
  সেটি টুরিং-গণনসাধ্য \emph{হয়}।
\end{prob}

\end{document}
